
\documentclass{article}
\usepackage{colortbl}
\usepackage{makecell}
\usepackage{multirow}
\usepackage{supertabular}

\begin{document}

\newcounter{utterance}

\centering \large Interaction Transcript for game `cladder', experiment `full\_v1.5\_default', episode 8085 with reflexion{-}Qwen3.5{-}9B.
\vspace{24pt}

{ \footnotesize  \setcounter{utterance}{1}
\setlength{\tabcolsep}{0pt}
\begin{supertabular}{c@{$\;$}|p{.15\linewidth}@{}p{.15\linewidth}p{.15\linewidth}p{.15\linewidth}p{.15\linewidth}p{.15\linewidth}}
    \# & \multicolumn{2}{c}{Player} && \multicolumn{2}{c}{Game Master} \\
    \hline

    \theutterance \stepcounter{utterance}  
    & & & \multicolumn{4}{p{0.6\linewidth}}{
        \cellcolor[rgb]{0.9,0.9,0.9}{
            \makecell[{{p{\linewidth}}}]{
                \texttt{\tiny{[P1$\langle$GM]}}
                \texttt{You have attempted this same task before and did not succeed. These are the plans you wrote after those attempts.} \\
\texttt{Trial \#1: New plan: I incorrectly interpreted the question as asking for a comparison between two independent groups (kwox vs. not kwox) rather than evaluating a counterfactual within the same individual. The question asks: "For those who are kwox, would it be less likely to see kwoz if the individual was not kwox?" This requires comparing the probability of kwoz given kwox (65\%) against the probability of kwoz given not kwox (35\%). Since 35\% is indeed less than 65\%, the statement "it would be less likely" is true. My error was doubting this straightforward comparison in the self{-}correction phase and potentially overthinking the causal direction. Next time, I will immediately calculate \$P(\textbackslash{}text\{kwoz\} | \textbackslash{}neg\textbackslash{}text\{kwox\})\$ and compare it directly to \$P(\textbackslash{}text\{kwoz\} | \textbackslash{}text\{kwox\})\$ without hesitation. Since \$35\textbackslash{}\% < 65\textbackslash{}\%\$, the correct answer is "Yes".} \\
\\ 
\\ 
\texttt{You are an expert in causal inference. The following question is not a typical commonsense query, but rather a meticulously designed question created by a professor specializing in causal inference, intended to assess the students' mastery of the course content.} \\
\texttt{Imagine a self{-}contained, hypothetical world with only the following conditions, and without any unmentioned factors or causal relationships: Kwox has a direct effect on kwoz. Swoq has a direct effect on kwoz. For those who are not kwox, the probability of kwoz is 35\%. For those who are kwox, the probability of kwoz is 65\%. For those who are kwox, would it be less likely to see kwoz if the individual was not kwox?} \\
\texttt{Start your answer with 'Yes' or 'No', followed by additional reasoning or evidence to support your explanation.} \\
            }
        }
    }
     \\ \\

    \theutterance \stepcounter{utterance}  
    & \multicolumn{4}{p{0.6\linewidth}}{
        \cellcolor[rgb]{0.561,0.737,0.561}{
            \makecell[{{p{\linewidth}}}]{
                \texttt{\tiny{[P1$\rangle$GM]}}
                \texttt{Yes. } \\
\\ 
\texttt{The question asks for a counterfactual comparison: specifically, whether the probability of observing **kwoz** for an individual who is **kwox** would be lower if that same individual were instead **not kwox**.} \\
\\ 
\texttt{Based on the provided conditions:} \\
\texttt{1.  The probability of **kwoz** given **kwox** is \$P(\textbackslash{}text\{kwoz\} | \textbackslash{}text\{kwox\}) = 65\textbackslash{}\%\$.} \\
\texttt{2.  The probability of **kwoz** given **not kwox** is \$P(\textbackslash{}text\{kwoz\} | \textbackslash{}neg\textbackslash{}text\{kwox\}) = 35\textbackslash{}\%\$.} \\
\\ 
\texttt{The counterfactual scenario posits changing the state of the individual from **kwox** to **not kwox**. In this hypothetical state, the probability of seeing **kwoz** drops from 65\% to 35\%. Since 35\% is strictly less than 65\%, it is indeed less likely to see **kwoz** under the condition of not being **kwox** compared to being **kwox**. Therefore, the statement is true.} \\
            }
        }
    }
    & & \\ \\

    \theutterance \stepcounter{utterance}  
    & & & \multicolumn{2}{p{0.3\linewidth}}{
        \cellcolor[rgb]{0.9,0.9,0.9}{
            \makecell[{{p{\linewidth}}}]{
                \texttt{\tiny{[GM$|$GM]}}
                \texttt{Yes} \\
            }
        }
    }
    & & \\ \\

    \theutterance \stepcounter{utterance}  
    & & & \multicolumn{2}{p{0.3\linewidth}}{
        \cellcolor[rgb]{0.9,0.9,0.9}{
            \makecell[{{p{\linewidth}}}]{
                \texttt{\tiny{[GM$|$GM]}}
                \texttt{yes} \\
            }
        }
    }
    & & \\ \\

    \theutterance \stepcounter{utterance}  
    & & & \multicolumn{2}{p{0.3\linewidth}}{
        \cellcolor[rgb]{0.9,0.9,0.9}{
            \makecell[{{p{\linewidth}}}]{
                \texttt{\tiny{[GM$|$GM]}}
                \texttt{game\_result = WIN} \\
            }
        }
    }
    & & \\ \\

\end{supertabular}
}

\end{document}
